% Generated draft: requires mathematical and source review.
\documentclass[11pt]{article}
\usepackage[margin=25mm]{geometry}
\usepackage{fontspec}
\setmainfont{Latin Modern Roman}
\usepackage{amsmath,amssymb,mathtools,unicode-math}
\setmathfont{Latin Modern Math}
\usepackage{enumitem,needspace,xurl,hyperref}
\setlength{\parindent}{0pt}
\setlength{\parskip}{5pt}
\setlength{\emergencystretch}{3em}
\newcommand{\N}{\mathbb{N}}
\newcommand{\Z}{\mathbb{Z}}
\newcommand{\Q}{\mathbb{Q}}
\newcommand{\R}{\mathbb{R}}
\newcommand{\C}{\mathbb{C}}
\newcommand{\F}{\mathbb{F}}
\newcommand{\A}{\mathbb{A}}
\newcommand{\PP}{\mathbb{P}}
\newcommand{\E}{\mathbb{E}}
\newcommand{\eps}{\varepsilon}
\newcommand{\dd}{\,\mathrm{d}}
\begin{document}
\section*{UnsolvedMath ID: 800238}
\textbf{Source problem code:} catalogue entry
\subsection{Definitions and mathematical statement}
Let \(n\geq 1\) and let \(\Omega\Subset\C^n\) be a pseudoconvex domain with smooth boundary; here \(\Omega\Subset\C^n\) means that its closure is compact. The supplied source asks whether
\[
\bar\partial:\Lambda^\alpha(\overline\Omega)\longrightarrow\Lambda_{(0,1)}^\beta(\overline\Omega)
\]
has closed range when \(\alpha<\beta\). The notation \(\Lambda^s\) denotes a regularity space whose precise convention and permitted exponents are not specified in the supplied record. The subscript \((0,1)\) denotes forms \(f=\sum_{j=1}^n f_j\,\dd\bar z_j\), with coefficients in the corresponding scalar regularity space. Distributionally on \(\Omega\),
\[
\bar\partial u=\sum_{j=1}^n\frac{\partial u}{\partial\bar z_j}\,\dd\bar z_j.
\]
The displayed operator requires a domain specification: membership in \(\Lambda^\alpha\) alone does not generally ensure that its derivative belongs to \(\Lambda^\beta\) when \(\alpha<\beta\).

One precise possible interpretation, conditional on fixing the regularity spaces and their topologies, is to set
\[
X_\alpha=\Lambda^\alpha(\overline\Omega),\qquad
Y_\beta=\Lambda_{(0,1)}^\beta(\overline\Omega),
\]
\[
\mathcal D_{\alpha,\beta}(\Omega)
=\{u\in X_\alpha:\bar\partial u\in Y_\beta\},\qquad
\mathcal R_{\alpha,\beta}(\Omega)
=\{\bar\partial u:u\in\mathcal D_{\alpha,\beta}(\Omega)\},
\]
where derivatives are understood distributionally on \(\Omega\). Under this interpretation, the question is whether, for every such domain and every permitted pair \(\alpha<\beta\),
\[
\overline{\mathcal R_{\alpha,\beta}(\Omega)}^{\,Y_\beta}
=\mathcal R_{\alpha,\beta}(\Omega).
\]
This restricted-domain formulation is a proposed clarification, not an identification verified from the source. The supplied record does not determine the convention for \(\Lambda^s\), the admissible exponent range, or the intended operator domain, so it does not yet specify a unique closed-range problem.
\subsection{Short English statement}
The source asks whether \(\bar\partial\) has closed range between the indicated regularity spaces on every smoothly bounded, relatively compact pseudoconvex domain when \(\alpha<\beta\). A precise formulation requires clarification of the spaces, exponents, and operator domain. The registry labels the problem open; this status is supplied metadata, not an independently verified assessment.
\subsection{Sources}
\begin{itemize}
\item Local UnsolvedMath registry, permanent ID \(800238\), source code \texttt{catalogue entry}, title \emph{Several Complex Variables Problem List --- Problem 35}. The supplied record attributes the problem to Buzzard, Fornaess, and Krantz, with proposed year 1995, and labels it open.
\item Buzzard, Fornaess, and Krantz, \emph{The Several Complex Variables Problem List} (1995), Problem 35. Supplied source URL: \url{https://arxiv.org/abs/math/9509202}. The record reports access on 2026-10-16 and extraction by \texttt{arxiv\_latex}; these are supplied provenance details, not a claim of independent source access or verification.
\end{itemize}
\end{document}
